% @author Lluis Gesa (lluis.gesa@uab.cat)
%
% Enginyeria Del Software
% -----------------------------------------------------------------------------------



% definim quin estil general volem per el document
\documentclass[11pt]{article}

% indiquem quins paquets amb funcionalitat extra volem fer servir
\usepackage{a4wide, amsmath, tabularx, colortbl, fancyhdr, graphicx, lastpage, caption}
\usepackage[utf8]{inputenc}

% Indiquem que comencem a definir el document i els seus continguts.
\begin{document}

% Definim valors i propietats, creem macros que ens ajudaran a
% agilitzar coses repetides:

\newcommand{\NomProjecte}{Perdidos en la ignorancia}
\newcommand{\NomDocument}{Procedimientos de pruebas}
\newcommand{\IdentificadorDocument}{GG-SL-N-TP}
\newcommand{\VersioDocument}{v1.0}

% Definim una macro que cada cop que insertem \requeriment, creara una
% taula amb el contingut passat per paràmetre
\newcommand{\TPStep}[5]{
\begin{center}
	\begin{tabularx}{\linewidth}{|p{2.5cm}|X|}
	\hline
	\cellcolor[gray]{.8} \textbf{Step} & \cellcolor[gray]{.9} #1\\
	\hline
	\cellcolor[gray]{.8} \textbf{Description} & #2 \\
	\hline
	\cellcolor[gray]{.8} \textbf{Input} & #3\\
	\hline
	\cellcolor[gray]{.8} \textbf{Output} & #4 \\
	\hline
	\cellcolor[gray]{.8} \textbf{Comments} & #5 \\
	\hline
	\end{tabularx}
\end{center}
}

% Fent servir propietas del paquet fancyhdr per crear una capçalera
% que es mostrarà cada pàgina
\pagestyle{fancy}
\fancyhead[LO,LE]{\includegraphics[width = 3.5cm]{./imatges/uab-logo.png}}
\fancyhead[CO,CE]{\bf \NomProjecte \\ \vspace{\baselineskip}
  \NomDocument }
\fancyhead[RO,RE]{
{\bf Ref : \IdentificadorDocument} \\
Versió: \VersioDocument\\
Data: \today \\
Pag. \thepage ~of \pageref{LastPage}
}
\fancyfoot[LO,LE,CO,CE,RO,RE]{}
\textheight=19cm


%
% Creem la pàgina inicial :

% Insertem el nom del document en una taula centrada feng servir una
% font gran \huge

\begin{center}
\begin{tabular*}{15cm}{p{14.5cm}}
\\
\\
\\
\rowcolor[gray]{.9}\begin{center}\huge{\bf{\NomDocument}}\end{center}\\
\\
\end{tabular*}
\end{center}


%
% Insertem una taula que contindrà les abreviacions que hi hagi al
% nostre document
%
\small{%
\begin{center}
\begin{tabular}{|p{2cm}|p{4.75cm}|p{2cm}|p{4.75cm}|}
\hline
\multicolumn{4}{|>{\columncolor[gray]{.9}[2mm][2mm]}c|}{{\bf Abreviaturas}} \\
\hline
UAB & Universitat Autonoma de Barcelona & ES & Enginyeria del Software \\
\hline
& & & \\
\hline
\end{tabular}
\end{center}
}

%
% Insertem una taula amb el historial de canvis del document
%
\small{%
\begin{center}
\begin{tabular}{|p{2cm}|p{2cm}|p{7.5cm}|p{2cm}|}
\hline
\multicolumn{4}{|>{\columncolor[gray]{.9}[2mm][2mm]}c|}{{\bf Historial de Revisiones}} \\
\hline
{\bf Version} & {\bf Date} & {\bf Comments} & {\bf Autor} \\
\hline
1.0 & 27-Mayo-2019 & Sprint 4 &
G43-2-02 \\
\hline
\end{tabular}
\end{center}
}




% Insertem en una nova pàgina, els indexs : de continguts, de figures
% i de taules, les comandes \tableofcontents ,  \listoffigures,
% \listoftables són gestionades per els paquets fets servir, i per
% defecte escriuen els titols de la secció en angles, per canviar-los,
% hem de 'renovar les comandes' :

\renewcommand{\contentsname}{Índex de continguts}
\renewcommand{\listfigurename}{Índex de figures}
\renewcommand{\listtablename}{Índex de taules}

% Ara li diem al LaTeX que inserti la informació
\newpage
\tableofcontents
\listoffigures
\listoftables

% Insertem les seccions que volem
\newpage

\section{Introducció}\label{sec:intro}


En este documento se pretenden recopilar todos los procesos o pruebas necesarias para comprobar que todas las partes del programa funcionan correctamente, es decir, que se ejecutan siguiendo los pasos correctos y responden de la forma esperada. \\ Para ello, se generarán diversos casos de test con la intención de probar todas las posibles combinaciones de entradas/salidas de los múltiples componentes que forman la aplicación.
\section{TP-01 : Registre}\label{sect:TP1}

En aquest Test Procedure verificarem tots els escenaris possibles
relacionat amb el registre de nous perfils al nostre sistema,
certificarem que es factible registrar nous usuaris,així com verificar
els límits relacionats amb el tamany i format de usernames i passwords.


\TPStep{01}{Inicia del sistema}{El sistema arranca amb una base de
dades buida sense cap usuari, unicament un administrador de password \textbf{admin1}}{Sistema
arrancat correctament, sense errors}{}

\TPStep{02}{Registre Nou usuari : Username}{S’accedeix al registre del
sistema, quan es demana username, s’introdueix un nom d’usuari nou: \textit{new\_user}}{El sistema detecta que l’usuari no existeix i demana si es vol crear}{}

\TPStep{03}{Registre Nou usuari : Password}{S’introdueix un password
valid : \textit{CXX965}!}{El sistema valida el password i crea el
perfil:
\begin{itemize}
\item A la base de dades hi ha el perfil nou
\item  El sistema mostra un missatge de benvinguda
\end{itemize}
}{}

\TPStep{04}{Registre Nou usuari : Username repetit}{S’accedeix al
registre del sistema, quan es demana username, s’introdueix un nom
d’usuari nou: \textit{new\_user}}{El sistema detecta que l’usuari existeix
i mostra error per pantalla}{}


\TPStep{05}{Registre Nou usuari : Username invalid}{S’accedeix al
registre del sistema, quan es demana username, s’introdueix un nom
d’usuari nou: \textit{aa}}{El sistema detecta que el nom d'usuari no te
prous lletres i mostra missatge d'error}{}


\TPStep{06}{...}{...}{...}{...}
\TPStep{07}{...}{...}{...}{...}

\section{TP-02 : Inicio de sesión}\label{sect:TP1}

En el siguiente caso de test se muestran los posibles inputs y outputs al realizar el login del usuario.


\TPStep{01}{Inicia del sistema}{El sistema se inicia sin ningún usuario por defecto. Previamente se debe haber creado manualmente en la base de datos con el usuario:contraseña root:toor}{El sistema arranca sin problemas}{}

\TPStep{02}{Inicio de sesión válido con contraseña válida : Usuario}{Se accede al login del sistema y se pide los datos de usuario: \textit{new\_user}}{El sistema pide la contraseña}{}

\TPStep{03}{Inicio de sesión válido con contraseña válida : Contraseña: Contraseña}{Se introduce la contraseña válida : \textit{CXX965D!}}{El sistema valida al usuario:
\begin{itemize}
\item  Redirige al usuario a su panel principal

\end{itemize}
}{}

\TPStep{04}{Inicio de sesión inválido con usuario inválido : Usuario}{Se introduce un usuario que no existe en la base de datos : \textit{inventado!}}{
\begin{itemize}
\item  El sistema comprueba en la base de datos y determina que no existe
\item Se devuelve un mensaje al usuario con el error
\end{itemize}
}{}

\TPStep{06}{Inicio de sesión inválido con contraseña inválida : Contraseña}{Se introduce un usuario y una contraseña que no coincide con la cuenta}{
\begin{itemize}
\item  El sistema comprueba en la base de datos y determina que no coincide el login con la contraseña
\item Se devuelve un mensaje al usuario con el error
\end{itemize}
}{}

\TPStep{07}{Inicio de sesión inválido con contraseña inválida: Control de
repeticiones}{Volvemos al paso 6 y lo repetimos 3 veces}{
\begin{itemize}
\item  El sistema va contando los inicios incorrectos y determina que debe bloquear el usuario por precaución
\item Envía al usuario un mensaje de que su cuenta ha quedado bloqueada por actividad sospechosa
\item Se envía un mensaje al administrador para que tome las medidas necesarias
\end{itemize}
}{}

\TPStep{08}{Inicio de sesión inválido con usuario bloqueado: Control de
repeticiones}{Se accede al registro y se determina que la cuenta introducida por el usuario está bloqueada}{El sistema devuelve al usuario un mensaje de error mostrándole los motivos}{}


\section{TP-03 : Eliminar cuentas}\label{sect:TP1}

En el siguiente test se mostrarán las diversas pruebas aplicadas a la eliminación de cuentas de usuario.


\TPStep{01}{Eliminar cuenta de usuario válido}{Se busca y selecciona en la base de datos el usuario que desea eliminar}
{
\begin{itemize}
\item  El sistema recibe la selección y procede a buscar y eliminar las entradas de la base de datos correspondientes a este usuario
\item El sistema devuelve al administrador un mensaje de confirmación
\end{itemize}}{}

\TPStep{02}{Eliminar cuenta de usuario inválido}{Se busca y selecciona en la base de datos el usuario que desea eliminar}
{
\begin{itemize}
\item  El sistema recibe la selección y procede a buscar y eliminar las entradas de la base de datos correspondientes, pero no encuentra ninguna con dicho usuario
\item El sistema devuelve al administrador un mensaje con el pertinente error
\end{itemize}}{}

\TPStep{03}{Eliminar cuenta de usuario sin permisos suficientes}{Se busca y selecciona en la base de datos el usuario que desea eliminar}
{
\begin{itemize}
\item  El sistema recibe la selección y procede a buscar y eliminar las entradas de la base de datos correspondientes, pero al intentar borrarlo detecta que el usuario no tiene suficientes permisos
\item El sistema devuelve al administrador un mensaje con el pertinente error
\end{itemize}}{}
\section{TP-04 : Borrar partida}\label{sect:TP1}

En el siguiente test se mostrarán las diversas pruebas aplicadas a la eliminación de partidas de usuario.


\TPStep{01}{Eliminar partida de usuario válida}{Se busca y se selecciona la partida a borrar}
{
\begin{itemize}
\item  El sistema recibe la petición de eliminación de la partida y busca en la base de datos su existencia
\item El sistema encuentra la partida y procede a borrar todas sus entradas
\item El sistema devuelve al usuario un mensaje de confirmación
\end{itemize}}{}

\TPStep{02}{Eliminar partida de usuario inválida}{Se busca y se selecciona la partida a borrar}
{
\begin{itemize}
\item  El sistema recibe la petición de eliminación de la partida y busca en la base de datos su existencia
\item El sistema devuelve al usuario un mensaje con el pertinente error
\end{itemize}}{}
\section{TP-05 : Modificar contraseña}\label{sect:TP1}

En el siguiente test se mostrarán las diversas pruebas aplicadas a la modificación de la contraseña de usuario.


\TPStep{01}{Cambio de contraseña válido}{Se ingresa la contraseña actual por motivos de seguridad y se ingresa la nueva contraseña dos veces}
{
\begin{itemize}
\item El sistema comprueba que la contraseña antigua es correcta
\item El sistema comprueba que la contraseña nueva cumple con los requisitos mínimos
\item El sistema comprueba que la contraseña nueva coincide en ambos campos
\item El sistema cambia la contraseña antigua por la nueva
\item El sistema envía un mensaje de confirmación al usuario
\end{itemize}}{}

\TPStep{02}{Cambio de contraseña inválido : Contraseña antigua incorrecta}{Se ingresa la contraseña actual por motivos de seguridad y se ingresa la nueva contraseña dos veces}
{
\begin{itemize}
\item El sistema comprueba que la contraseña antigua es incorrecta
\item El sistema envía un mensaje de error al usuario
\end{itemize}}{}

\TPStep{03}{Cambio de contraseña inválido : Contraseña nueva no cumple con requisitos mínimos}{Se ingresa la contraseña actual por motivos de seguridad y se ingresa la nueva contraseña dos veces}
{
\begin{itemize}
\item El sistema comprueba que la contraseña antigua es correcta
\item El sistema comprueba que la contraseña nueva cumple con los requisitos mínimos y detecta el error
\item El sistema envía un mensaje de error al usuario
\end{itemize}}{}

\TPStep{04}{Cambio de contraseña inválido : Contraseña nueva no coincide}{Se ingresa la contraseña actual por motivos de seguridad y se ingresa la nueva contraseña dos veces}
{
\begin{itemize}
\item El sistema comprueba que la contraseña antigua es correcta
\item El sistema comprueba que la contraseña nueva cumple con los requisitos mínimos
\item El sistema comprueba que ambas contraseñas nuevas coinciden y detecta el error
\item El sistema envía un mensaje de error al usuario
\end{itemize}}{}
\section{TP-06 : Estados de los usuarios}\label{sect:TP1}

En el siguiente test se mostrarán las diversas pruebas aplicadas a ver si los usuarios están conectados o no


\TPStep{01}{Comprobación de estado válida}{Se loguearán en el sistema 10 usuarios de muestra y se comprobará cuántos usuarios hay conectados}
{
\begin{itemize}
\item El sistema comprueba todas las cuentas en línea y las contabiliza
\item El sistema devuelve el número correcto de usuarios
\item El número de usuarios devuelto por el sistema, coincide con el número de usuarios logueados previamente 
\end{itemize}}{}

\TPStep{02}{Comprobación de estado inválida}{Se loguearán en el sistema 10 usuarios de muestra y se comprobará cuántos usuarios hay conectados}
{
\begin{itemize}
\item El sistema comprueba todas las cuentas en línea y las contabiliza
\item El sistema devuelve el número de usuarios conectados
\item El número de usuarios conectados no coincide con el devuelto por el sistema y, por tanto, es erróneo
\end{itemize}}{}
\section{TP-07 : Guardado manual}\label{sect:TP1}

En el siguiente test se mostrarán las diversas pruebas aplicadas para ver si el guardado manual del usuario funciona correctamente


\TPStep{01}{Guardado manual válido}{Se comienza una partida y en un momento dado se decide guardarla manualmente}
{
\begin{itemize}
\item El sistema recibe la petición y guarda todos los datos necesarios de la partida, para posteriormente retomarla
\item Se retoma la partida y continua en el punto esperado y con las propiedades esperadas

\end{itemize}}{}

\TPStep{02}{Guardado manual inválido}{Se comienza una partida y en un momento dado se decide guardarla manualmente}
{
\begin{itemize}
\item El sistema recibe la petición y guarda todos los datos necesarios de la partida, para posteriormente retomarla
\item Se retoma la partida pero no comienza por el punto esperado ni con las propiedades de la partida esperadas
\item Se devuelve el error 
\end{itemize}}{}

% Finalitzem el document.
\end{document}
